\section{Hợp chất hữu cơ và hydrocarbon}

\subsection{Hợp chất hữu cơ}

Hợp chất hữu cơ là hợp chất của carbon, trừ một số hợp chất đơn giản của carbon được xếp vào nhóm vô cơ như \ce{CO}, \ce{CO2}, \ce{H2CO3}, muối carbonate và bicarbonate.

\begin{ghinho}{Phân biệt hai khái niệm}
\begin{itemize}
\item \textbf{Hợp chất hữu cơ}: hợp chất của carbon theo quy ước trên.
\item \textbf{Hydrocarbon}: hợp chất hữu cơ chỉ chứa hai nguyên tố carbon và hydrogen.
\end{itemize}
\end{ghinho}

Ví dụ:
\begin{itemize}
\item \ce{CH4}, \ce{C2H6}, \ce{C2H4}: hydrocarbon;
\item \ce{C2H5OH}, \ce{CH3COOH}: hữu cơ nhưng không phải hydrocarbon;
\item \ce{CO2}, \ce{CaCO3}: vô cơ.
\end{itemize}

\subsection{Công thức phân tử và công thức cấu tạo}

Công thức phân tử cho biết số nguyên tử mỗi nguyên tố trong một phân tử. Công thức cấu tạo còn cho biết cách các nguyên tử liên kết với nhau.

Một số alkane đơn giản:
\begin{center}
\begin{tabular}{ccl}
\toprule
Tên & Công thức phân tử & Công thức cấu tạo thu gọn \\
\midrule
methane & \ce{CH4} & \ce{CH4} \\
ethane & \ce{C2H6} & \ce{CH3-CH3} \\
propane & \ce{C3H8} & \ce{CH3-CH2-CH3} \\
butane & \ce{C4H10} & \ce{CH3-CH2-CH2-CH3} \\
\bottomrule
\end{tabular}
\end{center}

Riêng \ce{C4H10} có hai công thức cấu tạo khác nhau:
\[
\ce{CH3-CH2-CH2-CH3}\qquad\text{và}\qquad\ce{CH3-CH(CH3)-CH3}.
\]
Hai chất có cùng công thức phân tử nhưng khác cách liên kết giữa các nguyên tử.

Từ các chất trên có thể nhận ra công thức chung của alkane mạch hở:
\[
\ce{C_nH_{2n+2}}\qquad (n\ge1).
\]

\subsection{Alkane và alkene}

Alkane chỉ có liên kết đơn giữa các nguyên tử carbon. Alkene có một liên kết đôi carbon--carbon trong phân tử.

Ethylene có công thức:
\[
\ce{C2H4},\qquad \ce{CH2=CH2}.
\]

Liên kết đôi làm ethylene có một số phản ứng khác methane và các alkane đơn giản, đặc biệt là phản ứng cộng bromine.

\begin{vidu}
Trong các chất \ce{CH4}, \ce{C2H4}, \ce{C2H5OH}, \ce{CO2}, chất nào là hydrocarbon? Chất nào thuộc alkene?
\end{vidu}
\begin{loigiai}
\ce{CH4} và \ce{C2H4} chỉ chứa C, H nên là hydrocarbon. \ce{C2H4} có liên kết đôi \ce{C=C} nên thuộc alkene.

\textbf{Vậy,} hydrocarbon là \ce{CH4}, \ce{C2H4}; trong đó \ce{C2H4} là alkene.
\end{loigiai}

\begin{tuluyen}
Phân loại các chất sau thành: hydrocarbon, hợp chất hữu cơ không phải hydrocarbon, hợp chất vô cơ.
\begin{tasks}(2)
\task \ce{C3H8};
\task \ce{CH3COOH};
\task \ce{Na2CO3};
\task \ce{C2H5OH};
\task \ce{C2H4};
\task \ce{CO2}.
\end{tasks}
\dapso{Hydrocarbon: \ce{C3H8}, \ce{C2H4}; hữu cơ không phải hydrocarbon: \ce{CH3COOH}, \ce{C2H5OH}; vô cơ: \ce{Na2CO3}, \ce{CO2}.}
\end{tuluyen}
